\chapter{Requirements Analysis and Specification}

\section{Introduction}

This chapter presents a comprehensive analysis of the requirements for the development
of the Sahhty platform, a mobile health application dedicated to pregnancy monitoring,
personalized medical care, and doctor-patient coordination.

Requirements analysis is a critical phase in the software development lifecycle. It defines
the expected functionalities, system constraints, and user interactions before any
implementation begins, ensuring that the delivered system faithfully addresses the real
needs of its stakeholders.

The chapter first identifies the system actors and their roles, then presents
the functional requirements organized across nine feature areas: authentication
and account management, patient health monitoring, pregnancy tracking, alert
management, doctor discovery and appointment booking, medication management
and treatment, medical file management, doctor features, and profile and
settings. The non-functional requirements are subsequently presented across
five quality dimensions: performance, security, scalability, privacy and data
protection, and usability and compatibility. It then introduces the UML global
use case diagram, followed by detailed refinements of the patient and doctor
use case diagrams, each accompanied by a descriptive table clarifying
scenarios, preconditions, and main flows.


Finally, the chapter concludes with a set of wireframe mockups illustrating
the key screens of the Sahhty application, providing a visual representation
of the user interface before the implementation phase.

% -------------------------------------------------------------
\section{Actor Identification}

The Sahhty platform involves four actors: two primary user roles, one specialized patient
role, and one automated system actor. Each actor is assigned a distinct set of privileges
tailored to their specific objectives.

\begin{itemize}
    \item \textbf{User (Base Actor):} Represents the common base shared by both patients
    and doctors. The User actor groups the features accessible to all authenticated
    accounts regardless of role, including authentication, profile management, and
    alert management. Both the Patient and Doctor actors inherit from User through
    a generalization relationship.

    \item \textbf{Patient (Primary User):} Represents any registered individual, male or
        female, who uses the application to monitor their health. Patients interact with the
        system to record health measurements, book and manage appointments with doctors,
        consult the national medication database, manage their active treatment plans,
        upload and organize personal medical documents, control doctor access to their
        medical files, and receive AI-generated risk alerts based on their health data.

            (\textbf{Female Patient} is a specialization of the Patient actor that inherits all
    general patient functionalities and additionally accesses features exclusive to female
    users: pregnancy lifecycle management, menstrual cycle tracking, and pregnancy risk
    evaluation during medication consultation.)


    \item \textbf{Doctor (Medical Professional):} Represents verified healthcare providers
    registered on the platform. Doctors manage their weekly schedule, confirm or cancel
    appointment requests from patients, and access patient medical files upon receiving
    explicit authorization. The doctor interface uses a dedicated green visual theme to
    clearly differentiate it from the patient interface.

    \item \textbf{System (Automated Actor):} Represents the backend processing layer,
    which includes the AI risk assessment engine (XGBoost classifier), the APScheduler
    background task manager, the Firebase Cloud Messaging notification service, and the
    Django Channels WebSocket server. The system acts autonomously to generate
    alerts, send push notifications, schedule reminders, and broadcast real-time events
    without requiring manual user intervention.
\end{itemize}

% -------------------------------------------------------------
\section{Functional Requirements}

The functional requirements define the specific behaviors and capabilities that the Sahhty
platform must deliver to each of its actors. Each requirement is categorized by feature
area, assigned a priority level, and broken down into measurable sub-requirements.

% --- Authentication ------------------------------------------
% --- Authentication ------------------------------------------
\subsection{Authentication and Account Management}

This section covers all features related to user registration, identity verification,
secure login, password recovery, and two-factor authentication for both patients and doctors.

\textbf{FR-1: Patient Registration}

\textbf{Priority:} High (patient registration is the mandatory entry point to the platform; all other features are inaccessible without a verified account.)

\textbf{Description:} The platform enables new patients to create a personal account by
providing personal and medical information across a multi-step registration form.

\textbf{Detailed Requirements:}
\begin{itemize}
    \item FR-1.1 \; The system collects first name, last name, email address, phone number,
    date of birth, gender, and password during the first registration step.
    \item FR-1.2 \; The password field enforces standard security rules: minimum eight
    characters, at least one uppercase letter, one digit, and one special character.
    \item FR-1.3 \; The system sends a six-digit OTP code by email after submission of
    the first step, and requires its verification before activating the account.
    \item FR-1.4 \; The second registration step collects medical profile data: height,
    weight, blood type, chronic diseases, current medications, family doctor name, and
    allergies selected from a DCI-based search interface.
    \item FR-1.5 \; For female patients, a third registration step collects menstrual cycle
    information (status, start date, end date) and pregnancy information (test date,
    start date, estimated due date).
    \item FR-1.6 \; The system rejects any registration attempt using an email address
    already associated with an existing account.
\end{itemize}

\medskip

\textbf{FR-2: Doctor Registration}

\textbf{Priority:} High (healthcare providers are a primary actor of the platform; their registration is a prerequisite for the appointment and medical file-sharing features to be operational.)

\textbf{Description:} The platform enables medical professionals to register with a
two-step flow that includes both account creation and professional profile setup.

\textbf{Detailed Requirements:}
\begin{itemize}
    \item FR-2.1 \; The first step collects identical basic account fields as patient
    registration (name, email, password, phone, date of birth, gender).
    \item FR-2.2 \; The second step collects professional information: medical speciality
    (selected from a predefined list), city, address, years of experience, consultation
    price, and a professional biography.
    \item FR-2.3 \; Upon successful registration, the doctor account remains in a
    pending verification state until approved by the platform administrator.
    \item FR-2.4 \; The system displays a verification status badge on the doctor home
    screen when the account is still pending approval.
\end{itemize}

\medskip

\textbf{FR-3: Login and JWT Authentication}

\textbf{Priority:} High (authentication is the security foundation of the platform; it gates access to all protected features and ensures that sensitive health data remains accessible only to verified users.)

\textbf{Description:} The platform authenticates users using email and password credentials
and issues a JWT access and refresh token pair upon success.

\textbf{Detailed Requirements:}
\begin{itemize}
    \item FR-3.1 \; The system verifies the submitted email and password against stored
    credentials and returns a JWT access token and a refresh token upon success.
    \item FR-3.2 \; The JWT payload encodes the user role (patient or doctor), allowing
    the frontend to route the user to the correct interface immediately after login.
    \item FR-3.3 \; The system stores the access and refresh tokens in the device's secure
    storage (Android Keystore) to persist sessions across application restarts.
    \item FR-3.4 \; The system automatically refreshes the access token using the refresh
    token when the access token expires, without requiring the user to log in again.
\end{itemize}

\medskip

\textbf{FR-4: Password Recovery}

\textbf{Priority:} High (without a recovery mechanism, any credential loss results in permanent account inaccessibility, directly compromising user retention and platform reliability.)

\textbf{Description:} The platform provides a three-step password recovery flow allowing
users to regain access to their account through email verification.

\textbf{Detailed Requirements:}
\begin{itemize}
    \item FR-4.1 \; The user initiates the recovery flow by submitting their registered
    email address on the forgot password screen.
    \item FR-4.2 \; The system sends a password reset OTP code to the provided email
    address and requires its verification before allowing a password change.
    \item FR-4.3 \; Upon successful OTP verification, the user is directed to a screen
    where they can define and confirm a new password meeting the platform's
    security rules.
    \item FR-4.4 \; The reset OTP code expires after a fixed time window and becomes
    invalid immediately upon first use to prevent replay attacks.
\end{itemize}

\medskip

\textbf{FR-5: Two-Factor Authentication (2FA)}

\textbf{Priority:} Medium (2FA strengthens account security but remains an optional, user-controlled mechanism that does not block access for users who choose not to enable it.)

\textbf{Description:} The platform offers an optional two-factor authentication mechanism
that patients and doctors can enable or disable from their settings screen.

\textbf{Detailed Requirements:}
\begin{itemize}
    \item FR-5.1 \; When 2FA is enabled, the system prompts the user for a time-based
    verification code after successful password verification.
    \item FR-5.2 \; The system allows the user to enable or disable 2FA at any time
    from the application settings screen.
    \item FR-5.3 \; An invalid or expired 2FA code returns a clear error message without
    granting access to the application.
\end{itemize}

% --- Patient Health Monitoring -------------------------------
\subsection{Patient Health Monitoring}

This section covers features that allow patients to record and track their vital health
measurements, receive AI-generated risk assessments, and optionally sync data from
a paired Wear OS smartwatch.

\textbf{FR-6: Health Measurement Recording}

\textbf{Priority:} High (health measurement recording is the central data-collection mechanism of the platform; the AI risk assessment, alert generation, and clinical consultation features all depend on it.)

\textbf{Description:} The platform allows patients to manually record health measurements
of multiple types and visualize them over time.

\textbf{Detailed Requirements:}
\begin{itemize}
    \item FR-6.1 \; The system supports the following measurement types: blood pressure
    (systolic and diastolic), heart rate, body temperature, blood glucose, and weight.
    \item FR-6.2 \; The input form adapts dynamically based on the selected measurement
    type: blood pressure requires two numeric fields, all other types require
    one value and a unit.
    \item FR-6.3 \; The measurements list screen supports filtering by type and sorting
    by date (newest-first or oldest-first), with paginated results.
    \item FR-6.4 \; After each successful submission, the backend returns the measurement
    record along with \texttt{risk\_level} and \texttt{risk\_percentage} fields computed
    by the AI engine.
\end{itemize}

\medskip

\textbf{FR-7: AI-Based Pregnancy Risk Assessment}

\textbf{Priority:} High (this feature is the primary clinical differentiator of the platform; it transforms raw health data into actionable medical insight and constitutes the core justification for the platform's maternal health focus.)

\textbf{Description:} The platform evaluates the patient's health risk level after each
measurement submission using a trained XGBoost classification model.

\textbf{Detailed Requirements:}
\begin{itemize}
    \item FR-7.1 \; The system returns a risk level (\textit{low}, \textit{medium}, or
    \textit{high}) and a numerical risk percentage with every measurement response.
    \item FR-7.2 \; A medium risk level triggers a warning dialog in the
    Flutter application, styled in orange, displaying the risk percentage.
    \item FR-7.3 \; A high risk level triggers a critical alert dialog styled in
    red, displaying the risk percentage and recommending immediate medical consultation.
    \item FR-7.4 \; A high risk result additionally creates a persistent alert in
    the patient's alert feed and dispatches a push notification via FCM.
    \item FR-7.5 \; A low risk result produces no dialog and continues normally.
\end{itemize}

\medskip

\textbf{FR-8: Smartwatch Integration (Wear OS)}

\textbf{Priority:} Medium (smartwatch integration enables passive heart rate monitoring but is supplementary; the manual measurement entry flow addresses the same clinical need, and the platform remains fully functional without it.)

\textbf{Description:} The platform integrates with Wear OS smartwatches to enable
passive, automated heart rate monitoring without requiring manual user input.

\textbf{Detailed Requirements:}
\begin{itemize}
    \item FR-8.1 \; A companion Kotlin application installed on the paired Wear OS watch
    uses WorkManager to trigger a heart rate reading every 30 minutes.
    \item FR-8.2 \; The watch application transmits the reading to the paired phone
    via the Wearable Data Layer API (DataClient).
    \item FR-8.3 \; The Flutter phone application receives the data through a
    \texttt{MethodChannel} named \texttt{com.sahhty/wearable} and automatically
    forwards it to the backend measurement endpoint.
    \item FR-8.4 \; Smartwatch-originated measurements are stored with the context
    value smartwatch data to distinguish them from manually entered readings.
\end{itemize}

% --- Pregnancy Tracking --------------------------------------
\subsection{Pregnancy Tracking}

This section covers features exclusively available to female patients, including
pregnancy progression monitoring and gestational context display.

\textbf{FR-9: Pregnancy Dashboard}

\textbf{Priority:} High (the pregnancy dashboard is a core feature for the platform's primary target demographic; it provides the gestational context essential for interpreting risk assessment results accurately.)

\textbf{Description:} The platform provides female patients with a dedicated pregnancy
tracking view displaying progression metrics and contextual information.

\textbf{Detailed Requirements:}
\begin{itemize}
    \item FR-9.1 \; The home screen displays a pregnancy progress card showing the
    current gestational week and day, computed from the recorded pregnancy start date.
    \item FR-9.2 \; The pregnancy tracking screen displays the estimated due date,
    pregnancy start date, elapsed weeks and days, and a visual progression indicator.
    \item FR-9.3 \; A fruit illustration represents the approximate fetal size at each
    stage of development, updated weekly.
    \item FR-9.4 \; The pregnancy dashboard and all related UI elements are hidden
    for male patients, who see a gender-appropriate interface without any pregnancy
    or menstrual cycle content.
\end{itemize}

% --- Alerts --------------------------------------------------
\subsection{Alert Management}

This section covers the alert and notification system that keeps both patients and
doctors informed of critical health events, appointment updates, and system messages
in real time.

\textbf{FR-10: Patient Alerts}

\textbf{Priority:} High (the alert system is a critical communication channel ensuring patients are promptly informed of high-risk health events, appointment changes, and medical file access requests.)

\textbf{Description:} The platform maintains a persistent alert feed for each patient,
populated by backend-generated health and system notifications.

\textbf{Detailed Requirements:}
\begin{itemize}
    \item FR-10.1 \; The alerts screen displays alerts categorized by type: All,
    Measurement Risk, Appointment, and System.
    \item FR-10.2 \; Each alert card displays a title, message body, severity indicator
    (color-coded left border), and creation timestamp.
    \item FR-10.3 \; The alerts list supports pagination, allowing the patient to browse
    their full notification history.
    \item FR-10.4 \; The patient receives real-time in-app notifications via WebSocket
    when a new medical file access request is submitted by a doctor or when a new
    appointment event occurs.
\end{itemize}

% --- Doctor Discovery ----------------------------------------
\subsection{Doctor Discovery and Appointment Booking}

This section covers features that allow patients to search for doctors, view their
profiles and availability, and book and manage appointments throughout their lifecycle.

\textbf{FR-11: Doctor Search and Filtering}

\textbf{Priority:} High (doctor discovery is a prerequisite for the appointment booking flow; without the ability to locate and evaluate healthcare providers, patients cannot initiate any consultation.)

\textbf{Description:} The platform allows patients to search and discover doctors using
multiple filtering criteria.

\textbf{Detailed Requirements:}
\begin{itemize}
    \item FR-11.1 \; The search interface supports filtering by doctor name, medical
    speciality, gender, and city (from a predefined list of the 24 Tunisian governorates).
    \item FR-11.2 \; A list view presents doctor cards showing name, speciality badge,
    city, years of experience, and consultation price.
    \item FR-11.3 \; An interactive map view displays the geographic locations of
    doctors as markers on an OpenStreetMap-based map, allowing patients to find
    nearby specialists visually.
    \item FR-11.4 \; The system supports proximity-based search, enabling patients to
    identify and retrieve the nearest doctors based on the patient's current location
    and the registered location of healthcare providers.
\end{itemize}

\medskip

\textbf{FR-12: Appointment Booking and Management}

\textbf{Priority:} High (appointment booking is the primary operational interaction between patients and doctors; it is the central use case that connects both user roles and delivers the platform's core value.)

\textbf{Description:} The platform allows patients to book appointments with doctors
and track their status throughout the booking lifecycle.

\textbf{Detailed Requirements:}
\begin{itemize}
    \item FR-12.1 \; The doctor profile screen displays the doctor's weekly availability
    schedule and a horizontal date picker highlighting days with available slots.
    \item FR-12.2 \; After selecting a date and confirming the appointment, the system
    submits the booking to the backend and notifies both the patient and the doctor
    via WebSocket, FCM push notification, and email.
    \item FR-12.3 \; The patient appointments screen displays all upcoming and past
    appointments, each showing the doctor's name, speciality, date, time, and current
    status (Pending, Confirmed, Completed, or Cancelled).
    \item FR-12.4 \; The patient appointments screen supports filtering by status
    (All, Pending, Confirmed, Completed) and sorting by date (recent or oldest),
    with paginated results.
\end{itemize}

% --- Medications ---------------------------------------------
\subsection{Medication Management and Treatment}

This section covers features that allow patients to manage their active treatments,
search the national medication database, and consult drug safety information including
interaction warnings and pregnancy risk classifications.

\textbf{FR-13: Treatment Management}

\textbf{Priority:} High (allergy incompatibility detection and automated dose reminders have direct patient safety implications, actively contributing to the prevention of medication errors.)

\textbf{Description:} The platform allows patients to manage their active treatment plans,
add medications with dose schedules, and receive automated medication reminders.

\textbf{Detailed Requirements:}
\begin{itemize}
    \item FR-13.1 \; The treatments screen displays each active treatment with its
    medication name, DCI (active ingredient), dosage schedule, start and end dates.
    \item FR-13.2 \; The patient can add a new treatment by selecting a medication,
    setting dose times, start and end dates, and frequency.
    \item FR-13.3 \; The system displays an allergy incompatibility badge next to any
    treatment whose DCI interacts with a DCI to which the patient has a recorded allergy.
    \item FR-13.4 \; The system automatically sends a medication reminder notification
    via FCM and email one hour before each scheduled dose time.
\end{itemize}

\medskip

\textbf{FR-14: National Medication Database Search}

\textbf{Priority:} High (access to the medication database enables informed pharmacological decisions; the LECRAT-based pregnancy risk classification is particularly critical given the platform's maternal health focus.)

\textbf{Description:} The platform allows patients and doctors to search the national
medication database and assess pregnancy safety at the DCI level.

\textbf{Detailed Requirements:}
\begin{itemize}
    \item FR-14.1 \; The search interface allows querying medications by commercial
    name, with results showing the commercial name, DCI, laboratory, and pregnancy
    risk category based on LECRAT classification.
    \item FR-14.2 \; The medication detail view presents trimester-by-trimester risk
    information, all known DCI-to-DCI interaction warnings, and allergy alerts when
    applicable.
    \item FR-14.3 \; A comparison tool allows the patient to evaluate two medications
    side by side, displaying their respective risk profiles and interaction data.
    \item FR-14.4 \; When a searched medication contains a DCI to which the patient is
    allergic, the system displays an allergy alert immediately in the search results.
    \item FR-14.5 \; Pregnancy risk evaluation is performed only for female patients
    with an active pregnancy record, and is skipped automatically for all other users.
\end{itemize}

% --- Medical Files -------------------------------------------
\subsection{Medical File Management}

This section covers features that allow patients to upload and organize their personal
health documents, and doctors to request and access authorized patient medical files.

\textbf{FR-15: Patient Medical Files}

\textbf{Priority:} High (centralized document management supports continuity of care, while the consent-based sharing mechanism ensures medical records are exchanged securely and only with authorized practitioners.)

\textbf{Description:} The platform allows patients to upload, organize, view, and share
their personal health documents.

\textbf{Detailed Requirements:}
\begin{itemize}
    \item FR-15.1 \; Patients can upload documents (images and PDFs) organized
    by category: Lab Results, Imaging, Prescriptions, Consultations, and Vaccinations.
    \item FR-15.2 \; Tapping a file opens an in-app viewer displaying the image or PDF
    without leaving the application.
    \item FR-15.3 \; Patients can manage which doctors are authorized to view
    their files through a dedicated access management section, where they can accept
    or reject pending access requests from doctors.
    \item FR-15.4 \; Patients can revoke a previously granted access at any time,
    immediately removing the doctor's permission to view their medical documents.
\end{itemize}

\medskip

\textbf{FR-16: Doctor Medical File Access}

\textbf{Priority:} High (authorized access to patient documents enables informed clinical decision-making, and the explicit consent model is fundamental to upholding the privacy standards expected of a medical platform.)

\textbf{Description:} The platform allows doctors to request access to a patient's medical
file and view it upon authorization.

\textbf{Detailed Requirements:}
\begin{itemize}
    \item FR-16.1 \; Doctors can search for patients by name and submit an
    access request from the patient search screen.
    \item FR-16.2 \; The patient receives a real-time WebSocket notification upon a
    new access request, allowing immediate response from any screen.
    \item FR-16.3 \; Once access is granted, the doctor can view all documents
    in the patient's file along with current measurements and active treatments.
    \item FR-16.4 \; Authorized doctors appear in the patient's list of granted-access
    practitioners, and the patient can revoke this access at any time.
\end{itemize}

% --- Doctor Features -----------------------------------------
\subsection{Doctor Features}

This section covers features specific to the doctor role, including appointment
management, weekly schedule configuration, and the dedicated alert feed.

\textbf{FR-17: Doctor Appointment Management}

\textbf{Priority:} High (without the ability to confirm or cancel appointments, the patient-side booking flow has no operational counterpart and cannot function as intended.)

\textbf{Description:} The platform allows doctors to manage all appointment requests
from their dedicated dashboard.

\textbf{Detailed Requirements:}
\begin{itemize}
    \item FR-17.1 \; The doctor appointments screen lists all incoming and past
    appointments with patient names, requested dates, and current statuses.
    \item FR-17.2 \; The doctor can confirm or cancel pending appointments
    directly from the appointments screen.
    \item FR-17.3 \; Confirming an appointment triggers a real-time notification to the
    patient via FCM, WebSocket, and email.
    \item FR-17.4 \; The appointments list supports filtering by status (All, Pending,
    Confirmed, Completed, Cancelled) and pagination.
\end{itemize}

\medskip

\textbf{FR-18: Doctor Schedule Management}

\textbf{Priority:} High (the doctor's availability schedule is the data source driving the patient-facing booking interface; inaccurate or missing entries would directly compromise the integrity of the appointment flow.)

\textbf{Description:} The platform allows doctors to define their weekly availability,
which is reflected in the patient-facing appointment booking interface.

\textbf{Detailed Requirements:}
\begin{itemize}
    \item FR-18.1 \; For each day of the week, the doctor can set a start time,
    an end time, and a number of available appointment slots.
    \item FR-18.2 \; A leave toggle switch allows the doctor to mark a specific day as
    unavailable without deleting the schedule entry, using the \texttt{is\_available}
    field in the backend model.
    \item FR-18.3 \; Schedule changes are reflected immediately in the appointment
    date picker visible to patients.
\end{itemize}

\medskip

\textbf{FR-19: Doctor Alerts}

\textbf{Priority:} Medium (the alert feed consolidates practice-related notifications and improves operational awareness, but doctors can manage their activity independently through the appointments and patient screens.)

\textbf{Description:} The platform provides doctors with a dedicated alert feed
aggregating all notifications relevant to their practice.

\textbf{Detailed Requirements:}
\begin{itemize}
    \item FR-19.1 \; Doctor alerts are organized into three tabs: Appointments
    (new bookings and cancellations), Patients (medical file access requests and
    acceptances), and System (platform notifications).
    \item FR-19.2 \; Each tab supports independent pagination.
    \item FR-19.3 \; Alert severity is communicated through a color-coded left
    border on each card, using the same convention as the patient alert interface.
\end{itemize}

% --- Profile and Settings ------------------------------------
\subsection{Profile and Settings}

This section covers profile management and application preference settings available
to both patients and doctors.

\textbf{FR-20: Patient Profile}

\textbf{Priority:} High (the profile screen gives patients visibility into their recorded personal and medical information, which is necessary to verify data accuracy and maintain trust in the platform.)

\textbf{Description:} The platform provides each patient with a profile screen
summarizing personal and medical information.

\textbf{Detailed Requirements:}
\begin{itemize}
    \item FR-20.1 \; The profile screen presents information in three expandable
    sections: Personal Information, Medical Information, and Menstrual Cycle
    (for female patients only).
    \item FR-20.2 \; Sections are collapsed by default and expand on tap to reduce
    cognitive load.
\end{itemize}

\medskip

\textbf{FR-21: Settings}

\textbf{Priority:} Medium (the settings screen provides useful personalization options such as language selection and notification preferences, but these do not affect the availability or correctness of the platform's core medical features.)

\textbf{Description:} The platform provides a settings screen accessible to both patients
and doctors with account management and application preferences.

\textbf{Detailed Requirements:}
\begin{itemize}
    \item FR-21.1 \; The settings screen offers language selection (French or English),
    two-factor authentication toggle, notification preferences, and account deletion.
    \item FR-21.2 \; For patients, an additional entry provides direct access to the
    medical file management screen.
    \item FR-21.3 \; For doctors, the settings screen includes access to the profile
    summary and the schedule management screen.
    \item FR-21.4 \; The account deletion action is irreversible and is highlighted
    in red to warn the user of its permanent consequence.
\end{itemize}

% -------------------------------------------------------------
\section{Non-Functional Requirements}

Non-functional requirements define the quality attributes and operational constraints that
the Sahhty platform must satisfy beyond its core features.

\subsection{Performance}

This section defines the expected responsiveness and processing speed of the platform
under normal and real-time usage conditions.

\textbf{NFR-1: API Response Time}

\textbf{Priority:} High (in a medical application where users expect immediate feedback, consistently slow API responses would erode user confidence and compromise the overall usability of the platform.)

\textbf{Description:} The platform delivers a responsive user experience under standard
network conditions.

\textbf{Detailed Requirements:}
\begin{itemize}
    \item NFR-1.1 \; The Flutter frontend renders initial content within three seconds
    on a stable mobile connection.
    \item NFR-1.2 \; The backend API endpoints return responses within
    500 milliseconds for standard operations such as login, profile retrieval, and
    measurement listing.
    \item NFR-1.3 \; AI risk assessment results are returned in the same measurement
    creation response, without requiring a secondary API call.
\end{itemize}

\medskip

\textbf{NFR-2: Real-Time Responsiveness}

\textbf{Priority:} High (delayed delivery of health alerts or appointment notifications could result in missed critical events, which is unacceptable for a platform intended to support patient safety.)

\textbf{Description:} Features that rely on live data exchange operate with negligible
latency.

\textbf{Detailed Requirements:}
\begin{itemize}
    \item NFR-2.1 \; WebSocket notifications (medical file access requests, appointment
    events) reach the recipient within one second of the triggering action over the real-time messaging layer.
    \item NFR-2.2 \; FCM push notifications are dispatched within two seconds of the
    backend event that triggers them.
    \item NFR-2.3 \; Smartwatch heart rate readings are transmitted to the backend
    within ten seconds of acquisition by the Wear OS sensor.
\end{itemize}

\subsection{Security}

This section defines the mechanisms in place to protect user credentials, health data,
and API endpoints from unauthorized access and common attack patterns.

\textbf{NFR-3: Credential and Data Protection}

\textbf{Priority:} High (health data is among the most sensitive categories of personal information; any compromise of credentials or medical records would constitute a serious breach of user trust with significant legal and ethical consequences.)

\textbf{Description:} The platform safeguards user credentials and sensitive health data
at every stage.

\textbf{Detailed Requirements:}
\begin{itemize}
    \item NFR-3.1 \; Passwords are hashed using the PBKDF2 algorithm before being stored in the database. Raw passwords never appear in any API response or log entry.
    \item NFR-3.2 \; JWT access tokens are short-lived and stored exclusively in the Android Keystore through the device's secure storage manager.
    \item NFR-3.3 \; All API communication occurs over HTTPS in production to
    prevent interception of health data in transit.
    \item NFR-3.4 \; Email OTP verification codes expire after a fixed window and
    become invalid immediately upon first use, preventing replay attacks.
\end{itemize}

\medskip

\textbf{NFR-4: Role-Based Access Control}

\textbf{Priority:} High (strict role separation at the API level is a fundamental requirement for a multi-role medical platform; without it, patient data could be accessed by unauthorized parties, violating the system's core privacy commitments.)

\textbf{Description:} Every protected API endpoint enforces authentication and
role validation at the backend level.

\textbf{Detailed Requirements:}
\begin{itemize}
    \item NFR-4.1 \; Every protected endpoint verifies the presence and validity of a
    JWT token before executing any handler logic.
    \item NFR-4.2 \; Patient-specific endpoints (measurements, pregnancy, alerts) return
    HTTP 403 if accessed by a doctor account, and vice versa.
    \item NFR-4.3 \; Medical file access requires explicit patient authorization; a
    doctor without a granted access record receives HTTP 403 when attempting
    to read patient documents.
    \item NFR-4.4 \; WebSocket connections validate the JWT token during the
    handshake and reject unauthorized connections immediately.
\end{itemize}

\subsection{Scalability}

This section defines the architectural decisions that allow the platform to handle
a growing number of concurrent users and increasing data volumes over time.

\textbf{NFR-5: Concurrent User Handling}

\textbf{Priority:} Medium (the selected architecture supports user base growth without structural rework; however, the current deployment scale does not yet require advanced concurrency optimizations beyond the chosen stack.)

\textbf{Description:} The system architecture accommodates a growing user base
without requiring structural changes.

\textbf{Detailed Requirements:}
\begin{itemize}
    \item NFR-5.1 \; The application server handles concurrent HTTP requests efficiently using multiple worker processes.
    \item NFR-5.2 \; Concurrent WebSocket connections are handled by a dedicated server component using an in-memory channel layer.
    \item NFR-5.3 \; PostgreSQL serves as the primary data store, supporting
    transactional integrity across concurrent writes from multiple active users.
\end{itemize}

\medskip

\textbf{NFR-6: Database Optimization}

\textbf{Priority:} Medium (indexing and server-side pagination improve query efficiency as data volumes grow, but these optimizations can be refined incrementally without blocking the initial deployment.)

\textbf{Description:} The database layer is structured to support frequent read
operations and complex filtered queries efficiently.

\textbf{Detailed Requirements:}
\begin{itemize}
    \item NFR-6.1 \; Frequently queried fields such as patient identifier, measurement
    type, and appointment status carry database indexes.
    \item NFR-6.2 \; Doctor location coordinates are stored as structured latitude and
    longitude fields to support geographic proximity queries used in the map discovery
    feature.
    \item NFR-6.3 \; Paginated API responses use server-side pagination to avoid
    transmitting large datasets to the mobile client.
\end{itemize}

\subsection{Privacy and Data Protection}

This section defines the controls and policies in place to ensure that sensitive patient
health data is accessed, stored, and transmitted in a secure and privacy-respecting manner.

\textbf{NFR-7: Patient Data Privacy}

\textbf{Priority:} High (the protection of personal health information is both a legal obligation and an ethical imperative; any unauthorized disclosure of medical records would constitute a critical failure of the platform's core responsibilities.)

\textbf{Description:} The platform enforces strict controls over access to sensitive
health and medical file data.

\textbf{Detailed Requirements:}
\begin{itemize}
    \item NFR-7.1 \; Medical file documents are never accessible to a doctor without
    prior explicit authorization from the patient.
    \item NFR-7.2 \; Health measurements, pregnancy records, and menstrual cycle data
    are accessible only to the owning patient and authorized doctors.
    \item NFR-7.3 \; All environment variables including database credentials, JWT secret
    keys, FCM service account keys, and email API keys are stored in a protected
    \texttt{.env} file excluded from version control.
\end{itemize}

\subsection{Usability and Compatibility}

This section defines the requirements ensuring the application remains accessible,
consistent, and easy to use across supported Android devices and languages.

\textbf{NFR-8: Mobile Compatibility}

\textbf{Priority:} High (the platform is delivered exclusively as a mobile application; inconsistent behavior across supported Android devices would directly prevent users from accessing the service, making broad device compatibility non-negotiable.)

\textbf{Description:} The mobile application functions correctly and consistently on
supported Android devices.

\textbf{Detailed Requirements:}
\begin{itemize}
    \item NFR-8.1 \; The frontend application targets Android 8.0 (API 26) and above as the minimum supported version.
    \item NFR-8.2 \; All screens are implemented using responsive Flutter layouts
    that adapt to different screen sizes and pixel densities without layout overflow.
    \item NFR-8.3 \; The application displays consistently across both physical Android
    devices and the Android Virtual Device emulator.
\end{itemize}

\medskip

\textbf{NFR-9: Accessibility and Localization}

\textbf{Priority:} Medium (supporting French and English broadens the platform's reach for a bilingual target audience; while localization meaningfully improves user experience, it does not affect the functional correctness of the platform's medical features.)

\textbf{Description:} The application is accessible and usable for patients regardless
of their preferred language.

\textbf{Detailed Requirements:}
\begin{itemize}
    \item NFR-9.1 \; The application supports French and English as fully localized
    languages, with all strings managed through the \texttt{flutter\_localizations}
    and \texttt{intl} packages.
    \item NFR-9.2 \; The active language is persisted across application restarts
    using SharedPreferences.
    \item NFR-9.3 \; All interactive elements meet minimum touch target size
    guidelines to remain usable for all patients, including those with limited
    dexterity.
\end{itemize}
% -------------------------------------------------------------
\section{Requirements Modeling}

\subsection{UML Modeling Language}

Unified Modeling Language (UML) is the internationally recognized standard for visually representing software systems. It provides a family of diagrams covering structure (class,
component, deployment) and behavior (use case, sequence, activity, state machine), enabling
developers, architects, and stakeholders to communicate about a system precisely and
unambiguously.

In the context of the Sahhty project, UML use case diagrams were produced at three levels of granularity: a \textbf{global diagram} encompassing all system actors and primary feature groups, a \textbf{patient-specific refinement} covering all patient interactions including female-specific flows, and a \textbf{doctor-specific refinement} detailing appointment management and patient file access workflows.

\begin{figure}[H]
    \centering
    \includegraphics[width=0.20\textwidth]{images/uml_logo.png}
    \caption{Unified Modeling Language (UML) logo}
    \label{fig:uml_logo}
\end{figure}

% -------------------------------------------------------------
\subsection{Global Use Case Diagram}

The global use case diagram (Figure~\ref{fig:uc_global}) presents the high-level
interactions between the four system actors (User, Patient, Doctor, and System) and the
main functional areas of the Sahhty platform. The User actor sits at the top of the
hierarchy and is connected to the three shared use cases inherited by both Patient and
Doctor through generalization relationships.

The Patient manages their health profile, records measurements, tracks pregnancy,
discovers doctors, books and manages appointments, consults medications, and manages
their medical files. The Doctor manages their weekly schedule, handles appointment
requests, and accesses authorized patient medical files. The System actor automatically
evaluates health risk after each measurement submission, dispatches scheduled
notifications, and sends automated alerts for missing measurements, unconfirmed
appointments, and pregnancy follow-up reminders.

\begin{figure}[H]
    \centering
    \makebox[\textwidth][c]{\includegraphics[width=1.3\textwidth]{images/uc_global.png}}
    \caption{Global Use Case Diagram of the Sahhty Platform}
    \label{fig:uc_global}
\end{figure}

\begin{tcolorbox}[colback=gray!10, colframe=gray!50, title=Access Control Rule]
All main features of the platform require prior user authentication.
Only the registration flow, OTP verification, and the login screen are publicly accessible.
Once authenticated, the system routes the user to the appropriate interface based on their
role (patient or doctor), ensuring that each actor can only access the features assigned
to their role.
\end{tcolorbox}

% -------------------------------------------------------------
\subsection{Patient Use Case Diagram}

The patient use case diagram (Figure~\ref{fig:uc_patient}) details all interactions
specific to the patient role. The diagram is organized into four sections: the User
section grouping inherited shared features, the Patient section covering all general
health and appointment features, the Female Patient section restricted to
pregnancy and menstrual cycle management, and the System section covering all
autonomous background tasks that interact with patient data.

Key relationships in this diagram include the \texttt{<<extend>>} relationship from
Verify 2FA to Authenticate, which fires only when the user has 2FA enabled, and the
\texttt{<<include>>} relationships from Consult Medication to Check Drug Interactions
and Check Allergy, which always execute on every medication consultation. The
pregnancy risk check extends Consult Medication conditionally for female patients
with an active pregnancy. Record Measurements includes Assess Health Risk when
all five measurement types are available, which in turn triggers alert generation
when a high or medium risk level is detected.

\begin{figure}[H]
    \centering
    \makebox[\textwidth][c]{\includegraphics[width=1.3\textwidth]{images/uc_patient.png}}
    \caption{Patient Use Case Diagram}
    \label{fig:uc_patient}
\end{figure}

Table~\ref{tab:uc_patient} provides the detailed textual description of the refined
patient use cases.

\begin{longtable}{|p{3.2cm}|p{1.6cm}|p{3.2cm}|p{5.4cm}|}
\caption{Detailed description of patient use cases}
\label{tab:uc_patient} \\
\hline
\textbf{Use Case} & \textbf{Actor} & \textbf{Preconditions} & \textbf{Main Scenario} \\
\hline
\endfirsthead
\hline
\textbf{Use Case} & \textbf{Actor} & \textbf{Preconditions} & \textbf{Main Scenario} \\
\hline
\endhead

Register Account &
Patient &
No existing account with this email &
1. Fill in personal information\newline
2. Verify email via OTP\newline
3. Complete medical profile\newline
4. (Female) Add pregnancy data\newline
5. Account activated \\
\hline

Login &
User &
Registered and verified account &
1. Enter email and password\newline
2. System validates credentials\newline
3. (If 2FA enabled) Enter verification code\newline
4. JWT token issued and stored\newline
5. Redirect to patient home screen \\
\hline

Forgot Password &
User &
Registered account &
1. Enter registered email address\newline
2. Receive OTP reset code by email\newline
3. Verify OTP code\newline
4. Set and confirm new password\newline
5. Redirect to login screen \\
\hline

View Home Dashboard &
Patient &
Authenticated &
1. Display pregnancy card\newline
2. Show latest measurements summary\newline
3. Show upcoming appointment\newline
4. Display recent alerts \\
\hline

Record Measurement &
Patient &
Authenticated &
1. Select measurement type\newline
2. Fill in value(s) and unit\newline
3. Submit to backend\newline
4. System returns risk level\newline
5. Display risk dialog if applicable \\
\hline

Assess Health Risk &
System &
All 5 measurement types recorded &
1. XGBoost model triggered automatically\newline
2. Risk level computed (low/medium/high)\newline
3. Risk note generated in French\newline
4. Alert created if risk is high or medium \\
\hline

View Measurements &
Patient &
At least one measurement recorded &
1. Load paginated measurements\newline
2. Apply type filter (optional)\newline
3. Toggle sort order (optional)\newline
4. Tap record to see details \\
\hline

Track Pregnancy &
Female Patient &
Active pregnancy record &
1. Open pregnancy screen\newline
2. View gestational week and day\newline
3. View due date and start date\newline
4. See fruit size analogy \\
\hline

Manage Menstrual Cycle &
Female Patient &
Authenticated female account &
1. Open menstrual cycle screen\newline
2. View cycle status and dates\newline
3. Log new period or end current period \\
\hline

View Alerts &
User &
Authenticated &
1. Load alerts by category tab\newline
2. Scroll through paginated list\newline
3. Tap alert to expand details\newline
4. Mark alert as read \\
\hline

Search Doctor &
Patient &
Authenticated &
1. Enter the doctor's name or select a medical specialty\newline
2. Apply optional filters (city, gender, or specialty)\newline
3. Optionally enable location-based search to display the nearest available doctors\newline
4. Browse results in list or map view\newline
5. Select a doctor card to view the complete profile \\
\hline

Book Appointment &
Patient &
Doctor profile open, slot available &
1. View doctor weekly schedule\newline
2. Select available date and time\newline
3. Confirm appointment booking\newline
4. System notifies doctor via WebSocket,\newline
\phantom{4. }FCM, and email \\
\hline

Cancel Appointment &
Patient &
Appointment in Pending or Confirmed status &
1. Open appointment details\newline
2. Tap cancel button\newline
3. System updates status to Cancelled\newline
4. Doctor receives cancellation notification \\
\hline

Consult Medication &
Patient &
Authenticated &
1. Search medication by name\newline
2. View results with risk category\newline
3. Open medication detail view\newline
4. System checks drug interactions\newline
5. System checks allergy conflicts\newline
6. (Female, active pregnancy) Display trimester risk \\
\hline

Manage Treatments &
Patient &
Authenticated &
1. View active treatments list\newline
2. Add medication with dose schedule\newline
3. System sets up medication reminders\newline
4. View allergy incompatibility warnings \\
\hline

Manage Medical Files &
Patient &
Authenticated &
1. View files grouped by category\newline
2. Upload new document\newline
3. Tap file to open in-app viewer\newline
4. Manage doctor access permissions \\
\hline

Grant Medical Access &
Patient &
Pending access request from doctor &
1. Receive real-time WebSocket notification\newline
2. Review request in access list\newline
3. Accept or reject the request\newline
4. System updates doctor permissions \\
\hline

Revoke Medical Access &
Patient &
Doctor has accepted access &
1. Open authorized doctors list\newline
2. Select doctor to revoke\newline
3. Confirm revocation\newline
4. Doctor access removed immediately \\
\hline

Update Settings &
User &
Authenticated &
1. Open settings screen\newline
2. Toggle 2FA, language, or notifications\newline
3. Changes saved immediately \\
\hline

\end{longtable}

% -------------------------------------------------------------
\subsection{Doctor Use Case Diagram}

The doctor use case diagram (Figure~\ref{fig:uc_doctor}) details all interactions
specific to the doctor role. The diagram is organized into three sections: the User
section grouping inherited shared features, the System section covering automated
appointment reminders and unconfirmed appointment alerts, and the Doctor section
covering all role-specific features.

The central use case of this diagram is Manage Appointment, which uses
\texttt{<<include>>} to always trigger Send Notification and View Appointment,
and uses \texttt{<<extend>>} relationships to optionally trigger Confirm Appointment
or Cancel Appointment based on the doctor's decision. The Request Medical Access
use case uses \texttt{<<include>>} from Search Patient, as the doctor must first
locate the patient before submitting an access request. View Patient Medical
Information extends into three specialized views (View Medical Files, View
Treatments, View Latest Measurements) available only after access is granted by
the patient.

\begin{figure}[H]
    \centering
    \makebox[\textwidth][c]{\includegraphics[width=1.3\textwidth]{images/uc_doctor.png}}
    \caption{Doctor Use Case Diagram}
    \label{fig:uc_doctor}
\end{figure}

Table~\ref{tab:uc_doctor} provides the detailed textual description of the refined
doctor use cases.

\begin{longtable}{|p{3.2cm}|p{1.6cm}|p{3.2cm}|p{5.4cm}|}
\caption{Detailed description of doctor use cases}
\label{tab:uc_doctor} \\
\hline
\textbf{Use Case} & \textbf{Actor} & \textbf{Preconditions} & \textbf{Main Scenario} \\
\hline
\endfirsthead
\hline
\textbf{Use Case} & \textbf{Actor} & \textbf{Preconditions} & \textbf{Main Scenario} \\
\hline
\endhead

Register Doctor Account &
Doctor &
No existing account with this email &
1. Fill in personal information\newline
2. Verify email via OTP\newline
3. Fill in professional profile\newline
4. Account set to pending verification \\
\hline

View Home Dashboard &
Doctor &
Authenticated and verified &
1. Display activity statistics\newline
2. Show pending and confirmed counts\newline
3. Show verification badge (if pending)\newline
4. Navigate to feature cards \\
\hline

Manage Appointment &
Doctor &
Authenticated &
1. View paginated appointment list\newline
2. Filter by status (optional)\newline
3. Select appointment to manage\newline
4. System loads appointment details\newline
5. Notification sent to patient \\
\hline

Confirm Appointment &
Doctor &
Appointment in Pending status &
1. Open pending appointment\newline
2. Tap confirm button\newline
3. Status transitions to Confirmed\newline
4. Patient notified via WebSocket,\newline
\phantom{4. }FCM, and email \\
\hline

Cancel Appointment &
Doctor &
Appointment in Pending or Confirmed status &
1. Open appointment details\newline
2. Tap cancel button\newline
3. Status transitions to Cancelled\newline
4. Patient receives cancellation notification \\
\hline

Manage Schedule &
Doctor &
Authenticated &
1. Open schedule management screen\newline
2. Set start time, end time, slots per day\newline
3. Toggle leave status if unavailable\newline
4. Save changes to backend \\
\hline

Search Patient &
Doctor &
Authenticated &
1. Enter patient name in search field\newline
2. View matching patient results\newline
3. Select patient to view profile\newline
4. Submit medical access request \\
\hline

Request Medical Access &
Doctor &
Patient found via search &
1. Select patient from search results\newline
2. Submit access request\newline
3. Patient receives real-time notification\newline
4. Status set to Pending approval \\
\hline

View Patient Medical Information &
Doctor &
Access granted by patient &
1. Open authorized patient entry\newline
2. Navigate to medical information tabs\newline
3. View medical files, treatments,\newline
\phantom{3. }or latest measurements \\
\hline

Update Location &
Doctor &
Authenticated &
1. Open profile settings\newline
2. Update city or GPS coordinates\newline
3. Location reflected on patient map view \\
\hline

\end{longtable}

% -------------------------------------------------------------

\section{Wireframe Mockups}

Wireframe mockups are low-fidelity visual representations of an application's user
interface, produced prior to the implementation phase to validate screen layouts,
navigation flows, and the organization of functional elements identified during
requirements analysis. In the context of the Sahhty project, six wireframes were
designed to cover the most critical screens of the application, spanning both the
patient and doctor interfaces.

% ── 2.6.1 ────────────────────────────────────────────────────
\subsection{Patient Home Dashboard}

The patient home dashboard constitutes the central entry point of the Sahhty
patient interface. Upon authentication, the patient is presented with a unified
view aggregating the most clinically relevant information: a pregnancy progression
card displaying the current gestational week and the corresponding fetal size
illustration, a summary strip of the latest recorded health measurements, the
next upcoming appointment with its status, and a preview of recent alerts. This
consolidated view was designed to minimize the number of navigation steps required
to access critical health data, directly supporting the usability requirements
defined in NFR-8 and NFR-9.

\begin{figure}[H]
    \centering
    \includegraphics[width=0.82\textwidth]{images/w01_dashboard.png}
    \caption{Wireframe mockup of the Patient Home Dashboard}
    \label{fig:wf_dashboard}
\end{figure}

% ── 2.6.2 ────────────────────────────────────────────────────
\subsection{Doctor Discovery and Appointment Booking}

The doctor discovery and appointment booking flow spans two consecutive screens.
The first screen presents the search interface with a specialty selector, city
filter, and gender filter, and allows the patient to switch between a list view
displaying doctor cards and an OpenStreetMap-based map view showing geographic
markers for nearby practitioners. Selecting a doctor opens the profile screen,
which displays the doctor's weekly availability schedule and a horizontal date
picker restricted to days with open slots. The patient selects an available time
slot and confirms the booking, after which the system simultaneously notifies the
doctor via WebSocket, FCM push notification, and email, as specified in FR-11
and FR-12.

\begin{figure}[H]
    \centering
    \includegraphics[width=\textwidth]{images/w02_doctor_discovery.png}
    \caption{Wireframe mockup of the Doctor Discovery and Appointment Booking flow}
    \label{fig:wf_booking}
\end{figure}

% ── 2.6.3 ────────────────────────────────────────────────────
\subsection{Health Measurement Entry and AI Risk Assessment}

The health measurement screen allows the patient to record a new vital sign
reading by selecting the measurement type and entering the corresponding values.
The input form adapts dynamically to the selected type: blood pressure exposes
two numeric fields for systolic and diastolic values, while all other types
require a single value and a unit selector. Upon submission, the backend
immediately returns the risk evaluation produced by the XGBoost classification
model. A medium risk result triggers an orange overlay dialog displaying the
computed risk percentage, while a high risk result triggers a red critical alert
dialog accompanied by a recommendation to seek immediate medical consultation
and a persistent entry in the patient's alert feed, as specified in FR-6 and
FR-7.

\begin{figure}[H]
    \centering
    \includegraphics[width=\textwidth]{images/w03_measurement_risk.png}
    \caption{Wireframe mockup of the Health Measurement Entry and AI Risk
    Assessment screens}
    \label{fig:wf_measurement}
\end{figure}

% ── 2.6.4 ────────────────────────────────────────────────────
\subsection{Medication Search and Drug Interaction Check}

The medication management screen provides patients and doctors with access to
the national medication database. The search interface returns result cards
annotated with the LECRAT pregnancy risk category badge and highlights any
allergy conflict detected against the patient's recorded allergy profile.
Opening a medication entry reveals the detail view, which presents a
trimester-by-trimester risk breakdown, the full list of DCI-to-DCI interaction
warnings, and a side-by-side comparison panel for evaluating two medications
simultaneously. For female patients with an active pregnancy record, the
pregnancy risk evaluation is computed and displayed automatically, while it is
suppressed for all other user profiles, as specified in FR-14.

\begin{figure}[H]
    \centering
    \includegraphics[width=\textwidth]{images/w04_medication.png}
    \caption{Wireframe mockup of the Medication Search and Drug Interaction
    Check screens}
    \label{fig:wf_medication}
\end{figure}

% ── 2.6.5 ────────────────────────────────────────────────────
\subsection{Medical File Management and Access Control}

The medical file management screen allows patients to upload, organize, and
consult their personal health documents, categorized into five types: Lab
Results, Imaging, Prescriptions, Consultations, and Vaccinations. Tapping any
file opens an in-app viewer without leaving the application. The access control
panel, accessible from the same screen, presents the list of doctors who have
submitted an access request alongside those who have already been granted
authorization. The patient can accept or reject any pending request and revoke
a previously granted access at any time, with the permission change taking
effect immediately on the backend, as specified in FR-15 and FR-16.

\begin{figure}[H]
    \centering
    \includegraphics[width=\textwidth]{images/w05_medical_files.png}
    \caption{Wireframe mockup of the Medical File Management and Access Control
    screens}
    \label{fig:wf_files}
\end{figure}

% ── 2.6.6 ────────────────────────────────────────────────────
\subsection{Doctor Appointment Management and Schedule Configuration}

The doctor interface consolidates appointment management and schedule
configuration into two dedicated screens. The appointment dashboard presents
all incoming and past requests in a paginated list with status filter tabs
covering Pending, Confirmed, Completed, and Cancelled states. Each appointment
card displays the patient's name, the requested date and time, and inline
confirm and cancel action controls, with every status change triggering a
real-time notification to the patient. The schedule configuration screen
allows the doctor to define weekly availability by setting a start time, end
time, and slot count for each day of the week, with a leave toggle switch
that marks a day as unavailable without removing the schedule entry from the
database, as specified in FR-17 and FR-18.

\begin{figure}[H]
    \centering
    \includegraphics[width=\textwidth]{images/w06_doctor.png}
    \caption{Wireframe mockup of the Doctor Appointment Management and Schedule
    Configuration screens}
    \label{fig:wf_doctor}
\end{figure}

\section{Conclusion}

This chapter provided a complete specification of the Sahhty platform's requirements,
covering all functional capabilities and quality constraints expected of the system.
Starting from a precise identification of the four system actors (User, Patient,
Female Patient, Doctor, and System), the chapter detailed twenty-one functional
requirements organized by feature area and nine non-functional requirements covering
performance, security, scalability, privacy, and usability.


